\chapter{Heuristics and biases: standards, environments and mechanisms}
\label{sec:p19-rationality-2-4}

A heuristic uses a tractable proxy or decision procedure in place of a
more demanding calculation or information search. Calling a procedure a
heuristic does not establish that it is inaccurate; calling an answer
biased does not identify the procedure that produced it. Fiedler, Prager
and McCaughey's 2021 treatment compares the Kahneman--Tversky
heuristics-and-biases program, the ecological fast-and-frugal program,
and the evidence needed to evaluate both. Its central methodological
question is whether a reported outcome diagnoses a cognitive mechanism,
an information environment, or their interaction
\cite[pp.~159--168]{FiedlerPragerMcCaughey2021Heuristics}.

The existing conjunction, natural-frequency and confidence studies at
\cref{sec:p19-conjunction-judgments,sec:p19-natural-frequencies-bayesian-reasoning,sec:p19-calibration-extreme-confidence}
are not three interchangeable measures of a general bias. They have
different response formats, event interpretations, samples and
benchmarks. This chapter retains the selected source's additional
families: availability, representativeness and small-sample expectations,
anchoring, counterfactual simulation, recognition, and ordered single-cue
choice. A lottery-choice process provides a further diagnostic comparison.
The evidence requirements are those of \cref{def:p19-evidence-record};
the bounded-rationality argument at \cref{sec:p17-rationality-8-5} and
the mental-model comparison at \cref{sec:p19-rationality-2-3} remain
distinct presentations.

\section{What is a bias relative to?}
\label{sec:p19-r24-standards}

\begin{hbdefinition}{def:p19-r24-benchmark}{A procedure, a target and an evaluation law}
Specify the question, the input information and its interpretation,
the target quantity or action, the evaluation law or reference record,
and the relevant loss and resource costs. A \emph{heuristic} is the
particular information-search, representation or choice procedure being
proposed. An \emph{error} is a deviation under the stated benchmark.
A statistical \emph{bias} is a systematic signed deviation under a
specified sampling law; a psychological bias label can instead name a
repeatable task pattern. Neither use licenses transferring that pattern
to every population or task.
\end{hbdefinition}

The source's perceptual analogy is the Ebbinghaus size contrast:
an equal-sized central circle can appear larger among smaller surrounding
circles and smaller among larger ones. Its argument is that an illusion
may expose a context in which a usually useful cue fails, rather than
prove that the whole processing strategy is generally inferior. That
analogy is not an estimate of how rare cognitive errors are. Their
frequency and practical cost require representative tasks and an
explicit evaluation law, and the claimed asymmetry between perceptual
and cognitive research remains an attributed methodological criticism
\cite[pp.~162,168, notes~1--2]{FiedlerPragerMcCaughey2021Heuristics}.

The source distinguishes \emph{coherence}, conformity to a stated
logical or probabilistic relation, from \emph{correspondence}, success
in an environment assessed through accuracy, speed and information cost.
Coherence requires stable meanings: a conjunction comparison cannot use
one interpretation for its parts and another for the whole.
Correspondence requires a target and criterion too: a bare prediction
accuracy such as seventy percent has no universal success threshold.
Costs of particular errors, a simple baseline, uncertainty and the
feasible alternatives affect its interpretation.

Unknown real-world parameters and limited computation can make a unique
optimal prescription unavailable. They do not erase every normative
constraint. A numerical answer can satisfy probability bounds yet use
poor evidence; a cheap decision can be useful while its stated
probability is incoherent. Neither observation proves that all formal
rules should be abandoned or that every normative method is feasible.
The source's call to study adaptation is a proposed research orientation,
not a theorem excluding normative evaluation
\cite[pp.~165--168]{FiedlerPragerMcCaughey2021Heuristics}.

\section{Availability: recall as a proxy and selective experience}
\label{sec:p19-r24-availability}

Availability uses assessed ease of retrieving or constructing relevant instances as a proxy for
frequency or probability. It can track frequent events when exposure,
retention and search preserve relevant frequency differences. It can
mislead when media coverage, recent experience, memorable features or
the search cue changes what is retrievable independently of the target.
The source uses judgments of causes of death and the first-versus-third
position of the letter K as examples. A memory store is not automatically
a representative random sample of a current population.

The source also cautions against choosing only a few exceptional letters
and generalizing their errors to the whole alphabet. The separate 1983
suffix-versus-position experiment at
\cref{sec:p19-conjunction-availability} already distinguishes a timed
production task from a corpus-frequency judgment. It is not the original
1973 K-letter task or the later larger-letter-set investigation.

\begin{hbconvention}{def:p19-r24-study-statistics}{Statistics in empirical study records}
In the empirical reports below, \(p\) denotes the source's reported
significance level or bound, and \(F\) and \(t\) denote its reported test
statistics. These are properties of the specified study and analysis,
not target-category probabilities or estimates of the probability that
a cognitive theory is true. The separately stipulated probability models
define their own \(p\), \(p_i\), and \(P\) explicitly.
\end{hbconvention}

The original 1973 Study~3 selected K, L, N, R and V because its reference
count put them more often in the third than the first position of English
words of at least three letters. Respondents judged which position was
more frequent and estimated the larger-to-smaller frequency ratio;
position order and three letter orders were varied. Among 152 respondents,
105 favored the first position for a majority of the letters, against 47
favoring the third; the reported sign test was \(p<.001\). Each letter's
median estimated ratio was 2:1 in the wrong direction. The article's
general recruitment record places Study~3 in a University of Oregon
student-newspaper recruitment pool, not a random population sample.
The selected letters and corpus definition restrict the result. Assessing
retrieval ease also need not mean explicitly retrieving examples on each
trial \cite[pp.~208--209,211--212]{TverskyKahneman1973Availability}.

Sedlmeier, Hertwig and Gigerenzer's 1998 judgment studies used
German words and first-versus-second positions, not the original
English first-versus-third comparison. Study~1 tested 114 paid
Salzburg students with four eight-letter lists, differing in
consonant/mixed composition and reversed order. Study~2 tested
390 Munich students individually, thirty per letter for thirteen
letters, each judging only one letter. Study~3 tested 77 paid
Salzburg students and preceded judgment with ninety-second
production tasks for O, U, N and R in the two positions.

The comparison distinguished retrieval amount, retrieval speed,
vowel/consonant class and regressed individual-frequency
predictions. The four accounts were compared by root-mean-square
deviation, contrast-analysis fit and consistency of choices with
predictions; the last accounted best for the stated estimate and
choice measures. Production-based predictions were also regressed
toward fifty percent for comparison. Producing examples beforehand had quantitative
but not the predicted qualitative effect. This constrains the
tested retrieval-proxy explanations, not every possible memory
account. Language, positions, selected letters, type/token
distinctions and the extra production task restrict generalization
\cite[pp.~759--762,765--767]{SedlmeierHertwigGigerenzer1998Letters}.

\begin{hbdefinition}{def:p19-r24-exposure}{A finite selection model for available experience}
Let \(p_i>0\), \(\sum_{i=1}^m p_i=1\), be the target proportions
of \(m\geq1\) event categories. Let \(s_i\in[0,1]\) be the
probability that a category-\(i\) event enters an exposure record, and
require \(\sum_i p_i s_i>0\). The record's conditional proportions
are denoted \(q_i\). This is a stipulated selection model, not a
measured model of any particular newspaper or memory system.
\end{hbdefinition}

\begin{hbtheorem}{thm:p19-r24-exposure}{proved=annex:p19.r24.exposure}{An exposure distortion without biased retrieval}
Under the selection model,
\[
 q_i=\frac{p_i s_i}{\sum_j p_j s_j}.
\]
Unbiased estimation of the record's proportions therefore targets
\(q_i\), not generally \(p_i\). The target ordering of two categories
can reverse in the record: for
\(p_1=.9,p_2=.1,s_1=.01,s_2=1\), the record has
\(q_1=9/109\) and \(q_2=100/109\).

For every probability vector \(q\) and every strictly positive target
vector \(p\) on this finite category set, some admissible selection
weights produce exactly \(q\). Thus observing a distorted retrieval
profile alone does not identify a distortion in retrieval as opposed
to prior exposure. This is a nonidentification result for the stipulated
model, not evidence that an actual exposure process has those weights.
\end{hbtheorem}

This contrast explains the source's criticism of attributing a
newspaper-related frequency error directly to an internal availability
mechanism. To separate the explanations, one needs an exposure record
and a retrieval manipulation or measurement, not just the final answer.
The relative usefulness of a proxy does not make its distortion harmless;
it locates what further information could correct the estimate
\cite[pp.~159--163]{FiedlerPragerMcCaughey2021Heuristics}.

Combs and Slovic's 1979 newspaper example concerns a content analysis,
not an experiment randomly assigning people to news exposure
\cite{CombsSlovic1979Newspaper}. Its relevance here is the distinction
between the frequency of a cause in an exposure record and its frequency
in the target population. The stipulated exposure construction proves
that this distinction can reverse an ordering; it neither estimates the
distortion in those newspapers nor identifies the cognitive process
producing a reader's judgment. Such empirical conclusions require the
study's sampling, coding, comparison measures and uncertainty.

\section{Representativeness: resemblance, base rates and samples}
\label{sec:p19-r24-representativeness}

Representativeness uses resemblance to a category or generating process
as a proxy for probability. The source's shy, withdrawn and helpful
person description invites comparison with a librarian stereotype rather
than a farmer stereotype. A diagnostic probability also depends on
category prevalence and on how frequently that description occurs in
each category. Without those quantities, selecting the stereotypical
profession does not by itself establish base-rate neglect.

For an explicitly stipulated two-category illustration, suppose only
one in 101 people is a librarian and the description is ten times as
likely for librarians as for farmers. Bayes' formula at
\cref{thm:p19-frequency-bayes-bridge} gives librarian posterior odds
\(10/100\), or posterior probability \(1/11\), rather than a
preference for the librarian category. The numerical premises are
invented benchmark data, not demographic estimates. The actual
natural-frequency task contrasts preserve their source-specific menus
and reference collections. The 1974 synopsis's experimental base-rate
comparison instead used descriptions allegedly sampled from one hundred
engineers and lawyers, with 70:30 or 30:70 category composition. These
are stated task pools, not respondent sample sizes or farmer/librarian
population estimates. For fixed nonzero description likelihoods, the
posterior-odds ratio between the two base-rate conditions is
\((.7/.3)^2=49/9\); the synopsis reports essentially unchanged
judgments rather than that shift, without supplying recruitment, group
sizes or confidence intervals for this comparison
\cite[pp.~1124--1125]{TverskyKahneman1974Heuristics}.

Representativeness also motivates the expectation that a small sample
should closely reproduce a population's proportions or that an irregular
exact sequence should be more likely than a regular-looking sequence.
The source's six-birth problem reports 72 families with exact sequence
GBGBBG and asks for an estimate of the number with BGBBBB. One sequence
contains three boys and three girls; the other five boys and one girl.
The corresponding coin illustration compares exact strings HTHTTH and
THTTTT. Exact strings must be distinguished from the larger class of
strings containing a specified count.

\begin{hbtheorem}{thm:p19-r24-strings}{proved=annex:p19.r24.strings}{Fair independent strings and the assumptions they need}
With six independent births and \(P(B)=p\), \(P(G)=1-p\),
the two stated birth strings have probabilities
\(p^3(1-p)^3\) and \(p^5(1-p)\). At \(p=1/2\), both
are \(1/64\). With independent coin tosses and head probability
\(p\), the two stated coin strings have probabilities
\(p^3(1-p)^3\) and \(p(1-p)^5\).
Fair marginal probabilities without independence do not guarantee
equal probabilities for these strings. Equal probabilities also do
not imply equal observed counts.
\end{hbtheorem}

\begin{hbtheorem}{thm:p19-r24-counts}{proved=annex:p19.r24.counts}{What the observed count of seventy-two does and does not fix}
Suppose \(N\) independently sampled six-birth families follow the
independent fair-birth model. Let \(X,Y\) count the two distinct
exact strings above. For \(N\geq72\),
\[
 E[Y\mid X=72]=\frac{N-72}{63}.
\]
This equals 72 when \(N=4608\); knowing \(X=72\) without knowing
\(N\) does not fix that conditional expectation. More generally,
for two disjoint categories of probabilities \(a,b\),
\(0\leq a<1\), \(a+b\leq1\), and any attainable count \(x\),
\[
 Y\mid X=x\ \sim\ \operatorname{Bin}
             \left(N-x,\frac{b}{1-a}\right).
\]
The case \(N=x\) is a degenerate count of zero. A plug-in estimate
using equal frequencies is a further estimation procedure, not a
deterministic equality imposed by the reported first count.
\end{hbtheorem}

The original 1972 task reports that 75 of 92 respondents estimated the
five-boy sequence as less frequent than the supplied three-boy sequence,
with median estimate 30 and sign-test \(p<.01\). The paper's general
methods describe classroom questionnaires, typically two to four
questions per respondent, in grades 10--12 of Israeli college-preparatory
schools; the approximately 1,500 respondents belong to the paper as a
whole, not this item. The item gives no total surveyed-family count or
confidence interval. A same-respondent order contrast also favored an
irregular three-boy string over BBBGGG, so ignoring order alone does not
explain every response. None of this is an observed demographic survey
of six-child families \cite[pp.~431--432]{KahnemanTversky1972Subjective}.

The experiment's interpretation concerns what a respondent assumes about
randomness and sampling. In an explicit fair independent-toss model,
irregularity is not an extra probability factor. In an unknown
real-world process, stationarity, fairness, dependence and sampling
may themselves need evidence. Removing an assumption changes the
benchmark; it does not retroactively make every error under the
original stated assumptions correct
\cite[pp.~160--161,166]{FiedlerPragerMcCaughey2021Heuristics}.

\section{Anchoring and counterfactual simulation}
\label{sec:p19-r24-anchoring-simulation}

Anchoring concerns a quantitative judgment's relation to an initial
number. The source's adjustment account proposes that a low or high
starting value is modified insufficiently as more information is
considered. Holiday budgeting illustrates a low starting price followed
by additions, or a high starting figure followed by deductions. A
starting number can instead affect accessibility of relevant information;
the outcome effect alone does not discriminate these mechanisms.

The selected source includes the wheel-of-fortune judgment about African
countries in the United Nations, auction starting prices and irrelevant
numerical primes. These are different tasks: an auction opening can
alter a market process as well as a bidder's numerical estimate. A
current claim about nations or prices is not inferred from the
historical task. Expertise likewise does not make an anchor irrelevant
by definition; its actual effect requires a specified sample and design.

Simulation uses ease of imagining a process or alternative outcome as
a cue for expectations and emotion. Missing a train by one minute can
make a successful counterfactual more accessible than missing it by
thirty minutes. The source's silver-versus-bronze Olympic illustration
compares possible upward and downward counterfactual reference points,
not a logical rule that every bronze medalist is happier than every
silver medalist. Rank, emotion, expectations and imagined alternatives
must be measured separately. This affective use of simulation is not
the same task as proving a railway procedure correct at
\cref{thm:p19-r23-reversal}
\cite[pp.~161--162]{FiedlerPragerMcCaughey2021Heuristics}.

The original 1974 synopsis reports a wheel-generated starting percentage,
a comparison with that number, and then an upward or downward estimate.
For the African-UN-membership item, the reported group medians were 25
and 45 after starting values of 10 and 65 respectively. That review does
not give the groups' sizes, population or a sampling-uncertainty interval
for these medians; they must not be supplied from the hypothetical
100-professional pool used in its separate engineer/lawyer task. The
given instruction to adjust also limits transfer to spontaneous
anchoring \cite[p.~1128]{TverskyKahneman1974Heuristics}.

Medvec, Madey and Gilovich's 1995 Study~1 used NBC's televised 1992
Barcelona Olympics footage: 23 silver and 18 bronze immediate reactions,
and 20 silver and 15 bronze medal-stand appearances. Twenty Cornell
undergraduates recruited for being uninformed and uninterested in sports
rated expressed emotion on a ten-point scale, with audio removed and
rank graphics obscured; five rated each of two versions of each tape.
The athlete-level means were 4.8 versus 7.1 for immediate silver/bronze
reactions and 4.3 versus 5.7 on the stand, with the reported medal main
effect \(F(1,72)=18.98,p<.001\). Seven athletes appeared in both
tapes; the authors acknowledge the initial independence problem and
report checks retaining only one appearance per athlete. They also
exclude events where bronze meant an immediate win and silver a loss,
and use a magazine's pre-event predictions as an expectations proxy.
These checks constrain alternatives but do not randomize medal rank,
equate a proxy with an athlete's actual expectations, or establish a
counterfactual mechanism uniquely
\cite[pp.~604--606]{MedvecMadeyGilovich1995Medalists}.

Study~2 had ten further Cornell raters assess interviews of thirteen
silver and nine bronze medalists. The upward-versus-``at least'' scale
differed as predicted, \(t(20)=2.37,p<.03\), but a second combined
category index did not reach a conventional .05 threshold
(\(p<.15\)). Study~3 directly questioned 115 silver/bronze swimming
and track medalists at the 1994 Empire State Games, 60 silver and
55 bronze, from scholastic and open divisions. Their ten-point thought
ratings averaged 6.8 versus 5.7, with SDs 2.2 and 2.7 and
\(t(113)=2.4,p<.02\). This is self-reported comparison direction,
not the first study's observer-rated happiness measure. The original
article provides no confidence intervals for these selected effects,
does not settle the contribution of social comparison, and explicitly
leaves persistence over time unresolved. It supports a contextual
association, not a recommendation to prefer a worse athletic result
\cite[pp.~606--609]{MedvecMadeyGilovich1995Medalists}.

\section{The lens model: validity and utilization}
\label{sec:p19-r24-lens}

The source's Figure~2.4.1 has a distal target, danger, and a subjective
danger judgment linked through five proximal cues: low familiarity or
recognition, availability of threatening events, simulation of such
events, depressing situational cues, and the base rate of dangerous
experience. The statistical relation between a cue and the distal target
is its \emph{ecological validity}; the weight or rule by which the
judgment uses that cue is \emph{cue utilization}. The diagram does not
make either relation a measured universal constant or a causal effect.
The metareasoning treatment distinguishes cue validity from cue use at
\cref{def:p19-r86-cue-standing}.

Depth perception supplies the source's analogy: binocular disparity,
motion, texture and blur can support distance judgments without a
specialized direct distance reading. \emph{Vicarious functioning}
allows one cue to replace another when conditions change. The source
also discusses lie detection from weak cues. A useful combination must
account for cue dependence, orientation, learned weights and the
evaluation environment; adding another weak cue does not, without
those assumptions, guarantee improvement. Nor does good cue utilization
repair an exposure record that no longer represents the target.
The distinction is between inference under uncertainty and unsupported
confidence, not between a fully direct perception and every imperfect
proxy \cite[pp.~162--163]{FiedlerPragerMcCaughey2021Heuristics}.

\section{Recognition and less-is-more effects}
\label{sec:p19-r24-recognition}

The source's pairwise recognition heuristic selects the recognized
object when exactly one of two objects is recognized and recognition
is positively related to the criterion. With both or neither recognized,
that particular cue does not discriminate. Knowledge that a town is
small or famous for an exceptional event can instead support a criterion
judgment or an evaluation of recognition's relevance. The chapter
contrasts the original noncompensatory recognition proposal with
Oppenheimer's objections and a later relevance-evaluation account.
These are different process specifications, not interchangeable uses
of the name recognition.

\begin{hbdefinition}{def:p19-r24-recognition-law}{A finite pair-comparison benchmark}
Take \(N\geq2\) objects with a strict target ranking, of which
\(n\) are recognized. Choose an unordered pair uniformly. On mixed
recognized/unrecognized pairs use recognition, whose conditional
correctness is \(\alpha\). On both-recognized pairs use a declared
knowledge-based procedure, whose conditional correctness is
\(\beta\). On both-unrecognized pairs guess fairly. Set a conditional
correctness parameter arbitrarily when its pair class is empty; it then
has zero weight. This is a stated decision model, not a premise that
human participants always obey it.
\end{hbdefinition}

\begin{hbtheorem}{thm:p19-r24-recognition-mixture}{proved=annex:p19.r24.recognition-mixture}{Pair classes and a qualified less-is-more condition}
The model's expected accuracy is
\[
 f(n)=\frac{2n(N-n)\alpha+n(n-1)\beta
                 +(N-n)(N-n-1)/2}{N(N-1)}.
\]
If the same \(\alpha,\beta\) are held fixed across recognition
amounts and \(\beta\geq1/2\), some \(n<N\) has
\(f(n)>f(N)\) exactly when \(\alpha>\beta\). In particular,
\(f(N-1)-f(N)=2(\alpha-\beta)/N\). For the original article's
stipulated \(N=100,\alpha=.8,\beta=.6\), recognizing fifty
objects gives \(6695/9900\), about 67.6\% accuracy, compared
with 60\% at complete recognition. The constant-validity assumption
does not imply that an actual person's newly recognized items leave
\(\alpha\) or \(\beta\) unchanged.
\end{hbtheorem}

Goldstein and Gigerenzer's 2002 study of heuristic accordance tested
22 University of Chicago students, native English speakers who had
lived in the United States for the preceding ten years. Six judged
all 300 pairs from the 25 largest German cities, and sixteen all
435 pairs from the thirty largest. Recognition was checked before
or after choice, with half in each order. Among pairs permitting
recognition use, mean accordance was 90\%, range 73--100\%.
Accordance is not answer accuracy and does not by itself exclude a
compensatory procedure producing the same choices
\cite[p.~81]{GoldsteinGigerenzer2002Recognition}.

A separate within-person study of 52 University of Chicago students,
native English speakers resident in the United States for the preceding
ten years, compared one hundred
random city pairs per country from the twenty-two largest German
and US cities, with country order counterbalanced. Mean answer
accuracy was 71.4\% for German cities and 71.1\% for US cities.
The report gives no uncertainty interval or significance test
establishing a population advantage from this .3-point difference.
A further study paid sixteen Munich residents to participate in four
weekly sessions. It introduced names from the seventy-five largest US
cities excluding the three largest, and used repeated population-choice
tests. In the fourth session, one hundred questions were drawn with
replacement from each person's first-session mixed-recognition pairs.
On these selected repeated questions, accuracy fell from 74.8\% to 71.3\%,
reported \(t(15)=1.82,p=.04\), one-tailed. That manipulation
is not equivalent to adding accurate criterion knowledge, and without
an unexposed control cohort it does not isolate every consequence of
repeated testing or newly acquired name familiarity
\cite[pp.~84--85]{GoldsteinGigerenzer2002Recognition}.

The source's broader claim is that ignorance can be useful in a
suitable environment. \Cref{thm:p19-r24-recognition-mixture} proves
one specified comparison; it does not recommend ignorance everywhere,
or identify a person's unobserved decision procedure. A reversal of
recognition's association with the criterion would change the useful
direction of the rule
\cite[pp.~164,167--168]{FiedlerPragerMcCaughey2021Heuristics}.

\section{Ordered single-cue choice and its computational evidence}
\label{sec:p19-r24-ttb}

Take the Best searches cues in an estimated validity order and stops
at the first discriminating cue, choosing the option favored by that
cue. Minimalist selects a cue order at random; Take the Last begins
with recently useful cues. Search, stopping and choice are distinct
parts of the procedure. Cue ordering is an input or learned property
to be explained, not a computation accomplished by naming it.
In particular, greater cue validity and greater discrimination rate
are different: a cue can be reliably correct when it discriminates
yet discriminate on few pairs.

The original 1996 binary TTB presentation also has a recognition
stage: choose the sole recognized object; guess if neither is
recognized; search the other cues if both are recognized. Cue values
can be positive, negative or unknown. Its published discrimination
rule stops when exactly one option has a positive value and the
other is negative \emph{or unknown}; a negative/unknown pair does
not discriminate. With no discriminating cue, it guesses. Unknown
is not an assertion of a negative fact, despite their sharing a
nonpositive role at that stopping test
\cite[pp.~652--654]{GigerenzerGoldstein1996FastFrugal}.

In the original article's Figure~1, two recognized objects with
positive/negative values on the first cue yield a choice of the
positive object immediately. For a different recognized pair, the
first cue's values are negative/unknown, so search continues; the
second cue's positive/negative comparison then decides. Lower-ranked
cues cannot compensate for either stopped choice. This is an execution
trace under a supplied ordering, not evidence that a person has learned
that ordering. The 1996 stopping test is also a version change from
the earlier positive/negative-only rule: positive/unknown now suffices
\cite[pp.~652--654, Figures~1--3]{GigerenzerGoldstein1996FastFrugal}.

That article's simulation used 83 German cities exceeding 100,000
inhabitants and nine ecological binary cues, with recognition supplied
as an additional cue to the integrative comparators. It combined
84 recognition amounts with six cue-knowledge amounts and generated
500 simulated individuals per combination. Each was evaluated on
all 3,403 city pairs; these are simulated states, not 252,000 human
participants. Recognition validity was set near .8 using a pilot
of 26 Chicago undergraduates. The competition included tallying,
weighted tallying, unit-weight and weighted linear models and
regression. Search cost counted accessed cue values, not measured
human response time \cite[pp.~654--658]{GigerenzerGoldstein1996FastFrugal}.

The original competition allowed regression to use the actual city
populations in fitting its weights and did not use a usual training/
held-out-test partition. Its result is a comparison in that specified
environment with missing knowledge, not an independent demonstration
of performance on a new data set. The authors also acknowledge that
their comparator variants are not every possible integration method.
Performance on held-out data requires a separate evaluation with a
specified training/test partition or independently sampled test cases;
the same-environment competition does not supply that evidence.
Simplicity can reduce fitting variance, but cannot guarantee superiority
in every environment; see the actual risk contrast at
\cref{thm:p17-br-bias-variance}. Likewise, binary cue dominance can
make lexicographic and weighted choice agree under the precise
inequalities in \cref{thm:p17-br-noncompensation}; a vague appeal
to correlated cues does not establish those inequalities.

\section{A diagnostic test of a lottery-choice procedure}
\label{sec:p19-r24-priority}

Brandst\"atter, Gigerenzer and Hertwig's original gain-domain
priority heuristic is a preference model, not a rule for inferring
the objectively larger item. For difficult, nondominating
two-outcome lotteries with nonnegative outcomes, it considers
minimum gain, probability of minimum gain, and maximum gain,
in that order. It stops at a sufficient minimum-gain difference
relative to the largest maximum, otherwise at a difference of
at least ten percentage points in minimum-gain probabilities,
otherwise chooses by the larger maximum. Gains are attractive
when larger, and the minimum-gain probability when smaller.
The source rounds one tenth of the largest maximum to a nearest
prominent number, drawn from powers of ten and their halves and
doubles. Ties in that rounding require an explicit convention;
the source's numerical illustrations avoid them
\cite[pp.~412--414]{BrandstaetterGigerenzerHertwig2006Priority}.

Its all-nonpositive loss variant instead prefers the smaller loss
and the larger probability of the smallest loss. With more than
two outcomes the original algorithm adds probability of the
maximum outcome as a fourth cue; it cannot be replaced by the
three-cue two-outcome version. Mixed-sign lotteries, dominance
and the excluded easy problems need their stated domain treatment,
not an assertion that the simple recipe covers all risky choice.
Predicting a majority preference is different from maximizing
expected money and from satisfying every expected-utility axiom.

\begin{hbexample}{def:p19-r24-priority-thresholds}{Rounding and inequality versions can change a prediction}
For a stipulated two-outcome gain pair, take one lottery paying 10 or
240, and another paying 32 or 200, each with probability one-half for
each outcome. Neither dominates; their expected monetary values are
125 and 116. The original nearest-prominent-number aspiration level
is 20, since one tenth of 240 is 24. The minimum-gain difference is
22, so that version stops and selects the 32-minimum lottery. An
unrounded threshold of 24 instead passes to equal minimum-gain
probabilities and then selects the lottery with maximum 240.

For another stipulated pair with equiprobable outcomes 0/200 and
20/180, both expected values are 100 and the aspiration level is 20
without a rounding difference. A difference of exactly 20 stops under
the original ``at least'' rule; a strict ``greater than'' rule proceeds
to the equal probabilities and then the maximum 200. These constructed
boundary cases compare precise algorithm versions; they are not
reported experimental items or evidence of a person's cognitive process.
\end{hbexample}

\begin{hblemma}{thm:p19-r24-priority-invariance}{proved=annex:p19.r24.priority-invariance}{A testable consequence of strict first-stage stopping}
Fix a declared outcome-dependent threshold and a pair of gain
lotteries whose minimum gains are unequal and whose difference meets its stopping
condition. A strictly first-stage-stopping procedure chooses by
that comparison regardless of changes in the lotteries'
probabilities that leave their outcomes and applicability domain
unchanged. A systematic probability influence in that region is
incompatible with this specified procedure. It does not identify
which alternative mechanism generated the influence.
\end{hblemma}

Fiedler's 2010 experiment tested fifty male and female students,
participating for pay or a study requirement, randomly assigned
to two forms of thirty-two hypothetical horse-betting choices.
Each form crossed minimum-outcome difference (one seventh or
one fourteenth of the largest maximum), minimum-outcome
probability difference (five or forty points), and baseline
winning probability (ten or thirty percent plus random increments),
with four tasks per cell. Choice was accompanied by pricing and
difficulty questions. The paper tests an \emph{unrounded}
one-tenth threshold specification, with a strict first-stage
inequality; it must not be silently identified with every rounded
boundary of the 2006 proposal.

Minimum-outcome probability affected choices even in the
tested first-stage region, with the reported probability-factor
\(F(1,49)=45.39,p<.001\). Overall accordance with the tested
priority predictions was .526, SD .103. In the one-seventh first-stage
region, mean risky choice was .875 with a forty-point probability
difference versus .611 with a five-point difference; the
outcome-by-probability interaction was \(F(1,49)=12.17,p<.01\).
The paper reports poor fit also on its subset of near-equal
expected-value lotteries, .529 for expected-value ratios at most 1.2.
No confidence interval for these selected statistics is supplied.
This is evidence against that strict process specification in
the administered task, not proof of irrational choices, universal
failure of heuristics, or superiority of every compensatory model
\cite[pp.~22--27]{Fiedler2010PriorityProcess}.

The simultaneous pricing and difficulty questions are relevant task
conditions, not process-neutral observations: price calculation could
encourage a different strategy. The paper discusses that objection,
but it does not make the choices diagnostic of every spontaneous
risky-choice situation. The test establishes a failure of a fixed
stopping specification for these materials; replacing it after observing
the results creates a new model requiring a new prediction test.

\section{Mechanism claims, expertise and the adaptive toolbox}
\label{sec:p19-r24-mechanisms}

An outcome predicted by both a fast procedure and a richer
procedure does not identify either one. A strong mechanism
comparison specifies which information is inspected, in what
order, when search stops, what knowledge is available, and
which response change is excluded by each procedure. The
availability example requires distinguishing exposure from
retrieval; the anchoring example requires distinguishing
adjustment from accessible information; a recognition account
requires distinguishing recognition from knowledge about its
source. A successful aggregate prediction can be useful without
answering all those process questions.

The source criticizes calling every mistaken response System~1
and every correct response System~2 after the result is known.
Such a labeling rule cannot independently predict which system
will be used. A constrained dual-process account may nevertheless
predict identifiable effects of time, learning or interference.
The possibility that systems overlap or that a single process
fits some findings is a competing architectural hypothesis,
not a proof that every two-process theory is logically impossible.
Likewise, numeracy or professional expertise is relevant information
about a sample, not a guarantee of perfect responding or proof
that an observed error came from unmotivated guessing.

The adaptive toolbox also needs an account of selecting a
procedure and assessing its environment. Adding another named
heuristic after every failed prediction supplies no such account.
Selection can have costs and use fallible cues; it need not
invoke an endless hierarchy if a specified learned or fixed
policy terminates the process. The source points to adaptive
strategy-selection and single-process alternatives while
emphasizing that decisive comparisons require constraints
\cite[pp.~166--168]{FiedlerPragerMcCaughey2021Heuristics}.

The chapter's final proposal emphasizes useful action under
uncertainty instead of treating a laboratory norm as a complete
account of real-world rationality. The handbook keeps the
benefit of that proposal without conflating distinct questions:
is the calculation coherent under its stated interpretation,
does its evidence describe the target, is its procedure feasible,
does it meet the stated decision loss, and what process do the
observations actually support? Uncertainty in one answer is
not permission to invent certainty in another.
